\documentclass[11pt,a4paper]{article}
\usepackage[margin=2.4cm]{geometry}
\usepackage{amsmath,amssymb,amsthm,mathtools}
\usepackage{booktabs}
\usepackage{enumitem}
\usepackage{xcolor}
\usepackage{hyperref}
\usepackage{fontspec}
\usepackage{polyglossia}
\setmainlanguage{russian}
\setotherlanguage{english}
\setmainfont{Times New Roman}
\setmonofont{Consolas}[Scale=0.90]

\hypersetup{
  colorlinks=true,
  linkcolor=blue!70!black,
  citecolor=blue!70!black,
  urlcolor=blue!70!black,
  pdftitle={Теоретическая физика дискретного планковского носителя, Часть III: Согласованная по импедансу бесфононная кремниевая фотовольтаика с делением синглетов (441 нм -> 2 x 882 нм < 1107 нм) и разрешение парадоксов интерфейса Тетрацен / Кремний в Lean 4},
  pdfauthor={Каратаев Дмитрий Александрович}
}

\theoremstyle{plain}
\newtheorem{theorem}{Теорема}[section]
\newtheorem{lemma}[theorem]{Лемма}
\newtheorem{proposition}[theorem]{Предложение}
\newtheorem{corollary}[theorem]{Следствие}
\theoremstyle{definition}
\newtheorem{definition}[theorem]{Определение}
\newtheorem{remark}[theorem]{Замечание}

\title{\textbf{Теоретическая физика дискретного планковского носителя, Часть~III:\\ Согласованная по импедансу бесфононная кремниевая фотовольтаика с делением синглетов ($441\;\text{нм}\to 2\times 882\;\text{нм} < 1\,107\;\text{нм}$) и разрешение парадоксов интерфейса «Тетрацен / Кремний» в Lean~4}}
\author{Д.\,А.~Каратаев}
\date{}

\begin{document}
\maketitle

\begin{abstract}
В третьей статье физической серии-сателлита к монографии \emph{«Теоретическая физика дискретного планковского носителя»} аппарат пошагового согласования фазовых и оболочечных импедансов на двусортной системе носителей $\langle\mathcal{M},\mathcal{T}\rangle$ (в смысле теории моделей А.\,И.~Мальцева; дискретный $M$-носитель $\mathbb{Z}_{512}$ на оболочке кубооктаэдра $O_2(3)$) применяется к задаче преодоления порога термализации Шокли--Квайссера ($33{,}7\%$) в кремниевой солнечной энергетике ($\mathrm{Si}$, $Z=14$, край поглощения $\lambda_{\mathrm{Si}}(3) = d^d\Delta\varphi_{\mathrm{disc}}(3) = 27\times 41 = 1\,107\;\text{нм}$, $E_g\approx 1{,}12\;\text{эВ}$). При несогласованном однопороговом поглощении синий фотон в главном планетарном окне $\lambda_{\mathrm{blue}}(3) = \Delta\varphi_{\mathrm{seam}}(3)^2 = 21^2 = 441\;\text{нм}$ ($2{,}81\;\text{эВ}$) теряет $\sim 60\%$ энергии ($\sim 1{,}69\;\text{эВ}$) на рождение надпороговых тепловых фононов ($\hat{n}_E > 0$), разогревая панель и снижая напряжение холостого хода~$V_{\mathrm{oc}}$. При делении синглетного экситона ($S_1\to T_1 + T_1$) удвоение длины волны рождает пару триплетов с $\lambda_{\mathrm{triplet}}(3) = 2\times 441 = 882\;\text{нм} = 6\times 147 = 18\times 49$ ($1{,}40\;\text{эВ}$), лежащих строго внутри окна поглощения кремния с квадратичным запасом $\Delta\lambda_{\mathrm{Si-SF}}(3) = 1\,107 - 882 = 225\;\text{нм} = 15^2 = (d\cdot d_{\mathrm{out}}(E))^2$.

Дано точное количественное разрешение трёх экспериментальных парадоксов гибридных интерфейсов «тетрацен / кремний» (2020--2026~гг.):
(1)~\emph{скачок $144\;\text{мэВ}$ и гашение синглетного деления на голом $\mathrm{Si}(100)$} (\texttt{arXiv:2310.15736}): полный заряд молекулы тетрацена $Z(\mathrm{C_{18}H_{12}}) = 18\times 6 + 12 = 120 = Q_\square(3)-|{\triangle}(3)|$ не достаёт $8$ единиц до четверть-кольцевого седла $Q_\square(3)=128$, а квадрат $12$-лучевой оболочки кубооктаэдра $|O_2(3)|^2 = 12^2 = 144$ в точности равен измеренному обменному сдвигу синглет-триплетной щели $144\;\text{мэВ}$;
(2)~\emph{кислородно-силанольная активация переноса триплетов на IBC-кремнии} (\texttt{arXiv:2007.06910}, \texttt{arXiv:2609.10979}): мостиковый атом кислорода $Z(\mathrm{O}) = 8 = |{\triangle}(3)|$ замыкает оболочку тетрацена ровно на седло ротатора $120 + 8 = 128 = Q_\square(3)$ ($\mathcal{C}_3(128)=0$, $\mathcal{S}_3(128)=+128$), а тетраэдрическая пассивация $\mathrm{SiO_4}$ или $[\mathrm{AlO_4H}]$ ($Z=46$) превышает порог затвора Паули $45^2 = 2\,025 < 2\,048 = K_P(3)^2/2 < 2\,116 = 46^2$;
(3)~\emph{подавление биэкситонной Оже-аннигиляции} (\texttt{arXiv:2606.31359}): разведение двух триплетов по ортогональным лучам кубооктаэдра при $\Delta\varphi = 128$ зануляет чётный интеграл перекрытия ($\mathcal{C}_3(128)=0$), каждый шаг переноса $\Delta\varphi_{\mathrm{seam}} = 21 < 41 = \Delta\varphi_{\mathrm{disc}}(3)$ идёт внутри мёртвой зоны Гейзенберга без рождения фононов ($\hat{n}_E(21)=0$), а согласование $n_1/n_2 = 9/8$ даёт просветление Френеля $T_{\mathrm{Fresnel}}(3) = 288/289\approx 99{,}654\%$. Все результаты формально верифицированы в \textbf{Lean~4} (\texttt{0 sorryAx}, \texttt{0 Classical.choice}).

\medskip\noindent
\textbf{Ключевые слова:} кремниевая фотовольтаика, деление синглетных экситонов, интерфейс тетрацен/кремний, предел Шокли--Квайссера, бесфононный перенос энергии, согласование импедансов на кубооктаэдре, Lean~4.
\end{abstract}

\section{Несогласованная фононная термализация и бесфононное деление синего кванта (\texorpdfstring{$441\;\text{нм}\to 2\times 882\;\text{нм} < 1\,107\;\text{нм}$}{441 nm -> 2 x 882 nm < 1107 nm})}
\label{sec:sf-wavelengths}

В двусортной структуре носителей $\langle\mathcal{M},\mathcal{T}\rangle$ (где термин \emph{носитель} / \emph{carrier} понимается в смысле теории моделей А.\,И.~Мальцева~\cite{Malcev1970}) кристаллический кремний ($\mathrm{Si}$, $Z = 14$) имеет край поглощения
\begin{equation}
  \label{eq:si-bandgap-1107}
  \lambda_{\mathrm{Si}}(3) \;=\; d^d\cdot\Delta\varphi_{\mathrm{disc}}(3) \;=\; 27\times 41 \;=\; 1\,107\;\text{нм} \qquad (E_g\approx 1{,}12\;\text{эВ}).
\end{equation}
При несогласованном поглощении синего фотона $\lambda_{\mathrm{blue}}(3) = \Delta\varphi_{\mathrm{seam}}(3)^2 = 21^2 = 441\;\text{нм}$ ($2{,}81\;\text{эВ}$) рождается лишь одна электрон-дырочная пара, а $\sim 1{,}69\;\text{эВ}$ ($60\%$ энергии) сбрасывается в решётку в виде тепловых фононов ($\hat{n}_E > 0$). При делении синглета ($S_1\to T_1 + T_1$) удвоение длины волны даёт два триплетных экситона на длине волны
\begin{equation}
  \label{eq:triplet-882}
  \lambda_{\mathrm{triplet}}(3) \;=\; \|e\|_2^2\,\lambda_{\mathrm{blue}}(3) \;=\; 2\times 441 \;=\; 882\;\text{нм} \;=\; |\square(3)|\cdot\mathcal{V}_{O_2}(3,3) \;=\; 6\times 147,
\end{equation}
лежащих строго внутри окна поглощения кремния~\eqref{eq:si-bandgap-1107} с точным квадратичным запасом:
\begin{equation}
  \label{eq:sf-margin-225}
  \Delta\lambda_{\mathrm{Si-SF}}(3) \;=\; \lambda_{\mathrm{Si}}(3) - \lambda_{\mathrm{triplet}}(3) \;=\; 1\,107 - 882 \;=\; 225\;\text{нм} \;=\; 15^2 \;=\; \bigl(d\cdot d_{\mathrm{out}}(E)\bigr)^2.
\end{equation}

\section{Разрешение трёх экспериментальных парадоксов интерфейсов «Тетрацен / Кремний» (2020--2026~гг.) и паспорт четырёхслойного каскада}
\label{sec:three-paradoxes}

\begin{theorem}[Согласование оболочечного и фазового импеданса в бесфононной фотовольтаике «Тетрацен--$\mathrm{SiO_x}$--Кремний»]
\label{thm:solar-boost}
При $d=3$ на кольце $\mathbb{Z}_{512}$ и оболочке $O_2(3)$ (\texttt{laplace\_demon\_impossibility\_and\_geometric\_maxwell\_demon\_operator} и \texttt{checkZeroPhononSingletFissionSiliconPhotovoltaicBoostOp} в \texttt{lean/MtCarrier/UnifiedOperators.lean}):
\begin{enumerate}[label=\textbf{(\roman*)},leftmargin=*]
\item \textbf{Оболочечный дефицит тетрацена на голом $\mathrm{Si}(100)$ и скачок $144 = 12^2\;\text{мэВ}$~\cite{Klymenko2023}:}
  заряд молекулы тетрацена $\mathrm{C_{18}H_{12}}$ равен $Z(\mathrm{C_{18}H_{12}}) = 18\times 6 + 12\times 1 = 120 = Q_\square(3) - |{\triangle}(3)| = 128 - 8$ (для пентацена $\mathrm{C_{22}H_{14}}$: $Z = 146 = \mathcal{V}_{O_2}(3,3)-1$). На голом кремнии ($Z=14$) оболочка тетрацена не достаёт $|{\triangle}(3)|=8$ единиц до седла расщепления $Q_\square(3)=128$, а квадрат $12$-лучевой оболочки кубооктаэдра $|O_2(3)|^2 = 12^2 = 144$ в точности равен измеренному сдвигу синглет-триплетной щели $144\;\text{мэВ}$.
\item \textbf{Кислородно-силанольный замок седла $120 + 8 = 128 = Q_\square(3)$ и барьер Паули $46^2 > 2\,048$~\cite{Daiber2020,McNab2026}:}
  введение на границе раздела атома кислорода $Z(\mathrm{O}) = 8 = |{\triangle}(3)|$ в составе тетраэдра кремнезёма $Z(\mathrm{SiO_4}) = 14 + 4\times 8 = 46$ (или водородно-пассивированного центра $Z([\mathrm{AlO_4H}]) = 13 + 32 + 1 = 46$):
  \begin{itemize}[noitemsep]
  \item замыкает оболочку тетрацена на седло ротатора $Z(\mathrm{C_{18}H_{12}}) + Z(\mathrm{O}) = 120 + 8 = 128 = Q_\square(3)$, где $\mathcal{C}_3(128) = 0$ (нулевая паразитная чётная связь) и $\mathcal{S}_3(128) = +128$ (максимум нечётного тока деления);
  \item помещает пассивирующий тетраэдр $Z=46$ непосредственно над порогом затвора Паули $K_P(3)^2/2 = 2\,048$:
    \[
      Z(\mathrm{AlO_4})^2 = 45^2 = 2\,025 \;<\; 2\,048 \;<\; 2\,116 = 46^2 = Z(\mathrm{SiO_4})^2 = Z([\mathrm{AlO_4H}])^2,
    \]
    блокируя преждевременный сброс неразделённого синглета в кремний и открывая резонансный перенос триплетов.
  \end{itemize}
\item \textbf{Подавление Оже-рекомбинации на ортогональных лучах, бесфононный перенос в зоне Гейзенберга и просветление Френеля $288/289$~\cite{Benning2026}:}
  при расхождении двух триплетов по ортогональным лучам кубооктаэдра с $\Delta\varphi = Q_\square(3) = 128$ их чётный интеграл перекрытия равен нулю ($\mathcal{C}_3(128) = 0$, подавление биэкситонного Оже-самотушения), каждый шаг переноса $\Delta\varphi_{\mathrm{seam}} = 21 < 41 = \Delta\varphi_{\mathrm{disc}}(3)$ излучает $\hat{n}_E(21) = 0$ тепловых фононов, а согласование $n_1/n_2 = \mathrm{mod\_bits}(3)/|{\triangle}(3)| = 9/8$ даёт коэффициент пропускания Френеля
  \[
    T_{\mathrm{Fresnel}}(3) \;=\; 1 - \left(\frac{9-8}{9+8}\right)^2 \;=\; \frac{288}{289} \;\approx\; 99{,}654\%.
  \]
\item \textbf{Замыкание каскада через репер планетарного дифференциала $\mathcal{F}_3$ и вронскиан Мейсона--Стотерса--Безу $\mathcal{W}_{\mathrm{MS}}$ на Замке~IV (\texttt{thm8}, \texttt{thm13} в \texttt{RotatorCalculusPaper.lean}):}
  весь четырёхслойный каскад кодируется кубическим этажом Замка~IV ($3\cdot 1^3 + 1\cdot 5^3 = 2\cdot 4^3 = 128$, валы $(X,Y,Z) = (1,5,4)$, $5^3 - 5 = 2^7 - 2^3 = 120 = Z(\mathrm{C_{18}H_{12}})$): в ортогональном репере дифференциала $\mathcal{F}_3$ ($\mathbf{r} = (1,5,4)\perp\mathbf{\omega} = (3,25,-32)$) квадрат нормы радиусов равен $\|\mathbf{r}\|_2^2 = 1^2 + 5^2 + 4^2 = 42 = S_2(3)$ (вторая сфера кубооктаэдра и ароматическое кольцо), проскальзывание шестерён равно $S_{AB} = 10 = Z(\mathrm{Ne})$, $M_{AB} = 40 = \Delta\varphi_{\mathrm{disc}}(3) - 1$ задаёт границу бесфононного коридора Гейзенберга, $M_{CA} = 44 = C_{\mathrm{III}} = Z(\mathrm{Si}) + Z(\mathrm{SiO_2})$ соответствует кварцево-кремниевому переходу Замка~III (чей параметр $k_{\mathrm{III}} = 15$ даёт запас поглощения $15^2 = 225\;\text{нм}$), след союзной матрицы равен $\operatorname{Tr}\mathcal{M}_3^{\mathrm{adj}} = 4^2 + 128 = 144 = 12^2\;\text{мэВ}$, а определитель четвёртых степеней целочисленного вронскиана Мейсона--Стотерса--Безу тождественно равен длине волны триплетного экситона:
  \begin{equation}
    \label{eq:wms-882}
    -\det\mathcal{W}_{\mathrm{MS}} \;=\; X^4 + Y^4 + Z^4 \;=\; 1^4 + 5^4 + 4^4 \;=\; 1 + 625 + 256 \;=\; 882\;\text{нм} \;=\; \lambda_{\mathrm{triplet}}(3).
  \end{equation}
\end{enumerate}
\end{theorem}

\begin{table}[htbp]
\centering
\caption{Четырёхслойный бесфононный каскад солнечной ячейки с согласованным импедансом ($d=3$).}
\label{tab:solar-cascade}
\small
\begin{tabular}{@{}llll@{}}
\toprule
\textbf{Слой каскада} & \textbf{Материал и геометрия} & \textbf{Инвариант носителя} & \textbf{Физическая функция} \\
\midrule
1. Входная оптика & Кварцевый микрорельеф $O_2(3)$ & $n_1/n_2 = 9/8$, $T = 288/289$ & Просветление $99{,}65\%$ и захват 12 лучей \\
2. Делитель кванта & Кристаллический тетрацен $\mathrm{C_{18}H_{12}}$ & $Z = 120$, $441\;\text{нм}\to 2\times 882\;\text{нм}$ & Деление синего фотона $S_1\to T_1+T_1$ ($\hat{n}_E=0$) \\
3. Седловой мост & Оксид/силанол $\mathrm{SiO_4}/[\mathrm{AlO_4H}]$ & $120+8=128=Q_\square$, $Z=46$ & Барьер Паули $46^2>2048$ для $S_1$, окно для $T_1$ \\
4. Базовый переход & Кремний IBC-$\mathrm{Si}$ (тыльные контакты) & $Z=14$, $\lambda_{\mathrm{Si}}=27\times 41=1107\;\text{нм}$ & Приём $2\times 882\;\text{нм}$ (запас $15^2=225\;\text{нм}$) \\
\bottomrule
\end{tabular}
\end{table}

\section{Формальная верификация в Lean~4}
Все теоремы формально верифицированы в \texttt{lean/MtCarrier/UnifiedOperators.lean} (\url{https://github.com/AI-Guiders/mt-lattice-fluid} и \url{https://github.com/AI-Guiders/scientific-papers}, \texttt{0 sorryAx}, \texttt{0 Classical.choice}).

\begin{thebibliography}{9}
\bibitem{Malcev1970}
А.\,И.~Мальцев, \emph{Алгебраические системы}, М.: Наука, 1970 (англ.\ пер.: Springer-Verlag, 1973).
\bibitem{Daiber2020}
B.~Daiber et~al., ``Singlet fission in tetracene/SiO$_x$/Si solar cells,'' \texttt{arXiv:2007.06910}, 2020.
\bibitem{Klymenko2023}
Y.~Klymenko et~al., ``Electronic structure and singlet fission of tetracene on Si(100),'' \texttt{arXiv:2310.15736}, 2023.
\bibitem{Benning2026}
M.~Benning, D.\,R.~Gamelin, F.\,T.~Rabouw et~al., ``Photon antibunching and Auger recombination in quantum-cutting phosphors,'' \texttt{arXiv:2606.31359}, 2026.
\bibitem{McNab2026}
S.~McNab, P.\,P.~Altermatt, T.\,W.~Schmidt, N.\,J.~Ekins-Daukes et~al., ``Singlet fission sensitized IBC silicon solar cells,'' \texttt{arXiv:2609.10979}, 2026.
\bibitem{Moura2021}
L.~de~Moura and S.~Ullrich, ``The Lean~4 theorem prover and programming language,'' in \emph{CADE~28}, LNCS vol.~12699, Springer, 2021, pp.~625--635.
\end{thebibliography}

\end{document}
